% Figure: unity-gain Sallen-Key second-order low-pass filter, drawn with plain
% TikZ (resistor and capacitor symbols hand-drawn, op-amp as a triangle). No
% circuitikz dependency.
\begin{figure}[htbp]
\centering
\begin{tikzpicture}[
  line width=0.8pt,
  res/.style={draw, engblue, fill=white, minimum width=10mm, minimum height=4mm, font=\scriptsize},
  font=\small,
]
  % input node
  \node[font=\scriptsize, anchor=east] at (-0.3,2.0) {$v_{\text{in}}$};
  \draw[engblue] (-0.3,2.0) -- (0.2,2.0);
  % R1
  \node[res] (R1) at (0.9,2.0) {$R_1$};
  \draw[engblue] (R1) -- (1.6,2.0);
  % node A (between R1 and R2)
  \node[circle, fill=engblack, inner sep=0.7pt] (A) at (1.6,2.0) {};
  \node[font=\tiny, anchor=south] at (1.6,2.1) {A};
  % R2
  \node[res] (R2) at (2.3,2.0) {$R_2$};
  \draw[engblue] (1.6,2.0) -- (R2) -- (3.0,2.0);
  % node B (op-amp non-inverting input)
  \node[circle, fill=engblack, inner sep=0.7pt] (B) at (3.0,2.0) {};
  \node[font=\tiny, anchor=south] at (3.0,2.1) {B};
  % C2 from node B to ground
  \draw[engblue] (3.0,2.0) -- (3.0,1.0);
  \draw[engblue, line width=1.1pt] (2.75,1.0) -- (3.25,1.0);
  \draw[engblue, line width=1.1pt] (2.75,0.85) -- (3.25,0.85);
  \node[font=\scriptsize, anchor=west] at (3.35,0.93) {$C_2$};
  \draw[engblue] (3.0,0.85) -- (3.0,0.5);
  \node[font=\tiny] at (3.0,0.3) {gnd};
  % op-amp triangle
  \draw[engblue] (3.4,1.4) -- (3.4,3.2) -- (5.2,2.3) -- cycle;
  \node[font=\scriptsize, anchor=west] at (3.5,2.8) {$-$};
  \node[font=\scriptsize, anchor=west] at (3.5,1.8) {$+$};
  % connect node B to the + input
  \draw[engblue] (3.0,2.0) -- (3.4,2.0);
  % output
  \draw[engblue] (5.2,2.3) -- (6.2,2.3);
  \node[font=\scriptsize, anchor=west] at (6.25,2.3) {$v_{\text{out}}$};
  % feedback wire from output to inverting input (unity gain)
  \draw[engblue] (5.7,2.3) -- (5.7,3.6) -- (3.2,3.6) -- (3.2,2.6) -- (3.4,2.6);
  % C1 from node A to the output node (positive feedback)
  \draw[engblue] (1.6,2.0) -- (1.6,3.0);
  \draw[engblue, line width=1.1pt] (1.35,3.0) -- (1.85,3.0);
  \draw[engblue, line width=1.1pt] (1.35,3.15) -- (1.85,3.15);
  \node[font=\scriptsize, anchor=east] at (1.3,3.07) {$C_1$};
  \draw[engblue] (1.6,3.15) -- (1.6,4.0) -- (5.7,4.0) -- (5.7,3.6);
  \node[font=\small] at (3.0,-0.1) {unity-gain Sallen--Key low-pass};
\end{tikzpicture}
\caption[Unity-gain Sallen--Key low-pass filter]{The unity-gain Sallen--Key
second-order low-pass filter. The op-amp is wired as a follower, so the
output equals the voltage at node B. The series resistors $R_1$ and $R_2$
roll the signal off, $C_2$ shunts node B to ground, and $C_1$ feeds the
output back to node A, the positive-feedback path that shapes the response
near the corner. The transfer function is second order, with a corner
frequency set by $R_1 R_2 C_1 C_2$ and a damping set by the resistor ratio
and the capacitor ratio, derived in the case study.}
\label{fig:vol04ch09:sallen-key}
\end{figure}
